\scanpage{250}

% head and opening, lines 1–5
\mspage{3}Ceci donnerait un principe de définition pour les opérations cherchées, Une difficulté vient du fait qu'en général, $\mathrm{Diag}$ n'est pas égal à (Cat), et que les types d'homotopie provenant de $\mathrm{Diag}$ ne \deleted{\unsure{sont}} \unsure{épuisent|forment|donnent} pas ($=$ à isom. près) \uline{tous} les types d'homotopie.

% upper middle, continuation and corrections, lines 6–11
\deleted{\unsure{Par} Mais on peut \unsure{se borner}} \inserted{Intuitivement, on} ne trouve que des types d'homotopie ``assez petits''. Mais on peut localiser $\mathrm{Diag}$ par ses \deleted{\unsure{qu}} \unsure{équ.|équi.} faibles (il \unsure{vaut|faut} mieux dire ``quasi-isomorphismes''), on trouve alors une catégorie \deleted{\unsure{plus}} $(\mathrm{Hot})_0$, et un foncteur canonique

% central arrow
\[\mathrm{Hot}_0\longrightarrow\mathrm{Hot}\ ,\]

% lower half, lines 1–7
si on a de la chance ce foncteur sera pleinement fidèle, et on \deleted{\unsure{défini}} essayera de décrire les foncteurs (1) ou (2) en y remplaçant Hot par la sous-catégorie pleine des ``petits'' types d'homotopie. On peut aussi oublier (provisoirement!) Hot, et travailler avec $\mathrm{Hot}_0$. C'est ce qu'on va faire pour le moment.

% before (5)
Pour pouvoir définir les foncteurs

% (5), product and Hom
\[
\left\{\begin{aligned}
\mathrm{Hot}_0^{\deleted{\unsure{\circ}}}\deleted{\unsure{\otimes}}\times\mathcal A&\longrightarrow\mathcal A&
(\unsure{\xi|\zeta}\unsure{\times|,} F)&\longmapsto\xi\times_{\mathbb D}F\\
\mathrm{Hot}_0^\circ\times\mathcal A&\unsure{-|\to}\mathcal A&
(\xi,F)&\longmapsto\mathrm{Hom}_{\deleted{\mathbb D}\mathbb D}(\xi,F)\ ,
\end{aligned}\right.\tag*{(5)}
\]

% foot, line carried over to next scan
je note que pour $F$ fixé, $X$ variable \deleted{\unsure{en}}dans $\mathrm{Diag}$, les

% archive sheet number: 250
